
\paragraph{\textbf{Corollary 2.4.}}
\SGARefTarget{sga3:xx:2.4}
\emph{Let}
\[
  \exp(Y)t\exp(X)
  \quad\text{and}\quad
  \exp(Y')t'\exp(X')
\]
\emph{be two elements of \(\Omega(S')\). Their product belongs to
\(\Omega(S')\) if and only if}
\[
  u=1+\langle X,Y'\rangle
\]
\emph{is invertible; in that case,}
\SGARefTarget{sga3:xx:2.4:F-prime}
\begin{equation}\tag{F'}
\begin{split}
  &\exp(Y)t\exp(X)\cdot\exp(Y')t'\exp(X')\\
  &\qquad
   =\exp\bigl(Y+u^{-1}\alpha(t)^{-1}Y'\bigr)\,
      tt'\alpha^\ast(u)\,
      \exp\bigl(u^{-1}\alpha(t')^{-1}X+X'\bigr).
\end{split}
\end{equation}

\paragraph{\textbf{Remark 2.5.}}
\SGARefTarget{sga3:xx:2.5}
\emph{Formula \textup{(F)} of \SGARefLink{sga3:xx:2.1}{Theorem 2.1} may also be written
without using the morphisms \(\exp\). Transporting the duality
\(\mathfrak g_\alpha\otimes\mathfrak g_{-\alpha}\to\mathscr O_S\)
through those morphisms gives a canonical pairing of vector
bundles}
\[
  U_\alpha\times_S U_{-\alpha}\longrightarrow\mathbb G_{a,S},
\]
\emph{which we continue to denote by
\((x,y)\mapsto\langle x,y\rangle\). Thus}
\[
  \langle\exp X,\exp Y\rangle=\langle X,Y\rangle.
\]
\emph{If \(x\in U_\alpha(S')\), \(y\in U_{-\alpha}(S')\), and
\(1+\langle x,y\rangle\in\mathbb G_m(S')\), then}
\[
  x\cdot y
   =y^{(1+\langle x,y\rangle)^{-1}}\,
      \alpha^\ast\bigl(1+\langle x,y\rangle\bigr)\,
      x^{(1+\langle x,y\rangle)^{-1}}.
\]

\paragraph{\textbf{Corollary 2.6.}}
\SGARefTarget{sga3:xx:2.6}
\emph{The pairing}
\[
  W(\mathfrak g_\alpha)\times_S W(\mathfrak g_{-\alpha})
    \longrightarrow\mathbb G_{a,S}
\]
\emph{induces a pairing of principal
\(\mathbb G_{m,S}\)-bundles}
\[
  W(\mathfrak g_\alpha)^\times
    \times_S W(\mathfrak g_{-\alpha})^\times
      \longrightarrow\mathbb G_{m,S}.
\]
\emph{We denote this pairing by
\((X,Y)\mapsto\langle X,Y\rangle\), or simply
\((X,Y)\mapsto XY\).}

For every \(X\in\Gamma(S,\mathfrak g_\alpha)^\times\), there is
therefore a unique
\[
  X^{-1}\in\Gamma(S,\mathfrak g_{-\alpha})^\times
\]
such that \(XX^{-1}=1\). Moreover,
\[
  (zX)^{-1}=z^{-1}X^{-1}.
\]
The morphism
\[
  s\colon W(\mathfrak g_\alpha)^\times
    \longrightarrow W(\mathfrak g_{-\alpha})^\times
\]
so defined is an isomorphism of schemes, compatible with the
isomorphism \(s\colon z\mapsto z^{-1}\) of their operator groups.

\paragraph{\textbf{Definition 2.6.1.}}
\SGARefTarget{sga3:xx:2.6.1}
\emph{The elements \(X\) and \(s(X)=X^{-1}\) are said to be
paired.}

Apply \SGARefLink{sga3:xx:2.4}{Corollary 2.4} with \(Y=0=X'\) and
\[
  Y'=aX^{-1},\qquad a\in\mathscr O_S(S).
\]
Then \(u=1+a\), and
\[
  u^{-1}Y'
   =u^{-1}(u-1)X^{-1}
   =(1-u^{-1})X^{-1}.
\]
Hence:

\paragraph{\textbf{Corollary 2.7.}}
\SGARefTarget{sga3:xx:2.7}
\emph{Let \(X\in\Gamma(S,\mathfrak g_\alpha)^\times\) and
\(u\in\Gamma(S,\mathscr O_S)^\times\). Then}
\[
  \alpha^\ast(u)
   =\exp\bigl((u^{-1}-1)X^{-1}\bigr)\,
      \exp(X)\,
      \exp\bigl((u-1)X^{-1}\bigr)\,
      \exp(-u^{-1}X).
\]

\paragraph{\textbf{Definition 2.8.}}
\SGARefTarget{sga3:xx:2.8}
\emph{The morphism \(\alpha^\ast\) is called the coroot associated
with the root \(\alpha\).}

\paragraph{\textbf{Remark 2.9.}}
\SGARefTarget{sga3:xx:2.9}
\emph{If \((G,T,\alpha)\) is an \(S\)-elementary system, then so
is \((G,T,-\alpha)\). \SGARefLink{sga3:xx:2.1}{Theorem 2.1} therefore gives a duality
between \(\mathfrak g_{-\alpha}\) and \(\mathfrak g_\alpha\), and
a coroot \((-\alpha)^\ast\). Taking the inverse of formula
\textup{(F)} gives immediately}
\[
  \langle X,Y\rangle=\langle Y,X\rangle,
  \qquad
  (-\alpha)^\ast=-\alpha^\ast.
\]

We now pass to the Lie algebra of \(G\). The root \(\alpha\) and
the coroot \(\alpha^\ast\) define linear forms
\[
  \mathscr O_S\xrightarrow{\ \overline{\alpha^\ast}\ }
    \mathfrak t\xrightarrow{\ \bar\alpha\ }\mathscr O_S.
\]
Put
\[
  H_\alpha=\overline{\alpha^\ast}(1).
\]
We call \(\bar\alpha\) the infinitesimal root associated with
\(\alpha\), and \(H_\alpha\) the corresponding infinitesimal
coroot.

\paragraph{\textbf{Lemma 2.10.}}
\SGARefTarget{sga3:xx:2.10}
\emph{Let \(S'\to S\), and let}
\[
  X,X'\in W(\mathfrak g_\alpha)(S'),\qquad
  H\in W(\mathfrak t)(S'),\qquad
  Y,Y'\in W(\mathfrak g_{-\alpha})(S'),\qquad
  t\in T(S').
\]
\emph{Then}
\[
\tag{1}\label{sga3:xx:lemma:2-10:equation:1}
  \operatorname{Ad}(t)H=H,\qquad
  \operatorname{Ad}(t)X=\alpha(t)X,\qquad
  \operatorname{Ad}(t)Y=\alpha(t)^{-1}Y;
\]
\[
  \tag{2}\label{sga3:xx:lemma:2-10:equation:2}
\begin{aligned}
  \operatorname{Ad}(\exp X)H
    &=H-\bar\alpha(H)X,&
  \operatorname{Ad}(\exp X)X'&=X',\\
  \operatorname{Ad}(\exp X)Y
    &=Y+\langle X,Y\rangle H_\alpha
        -\langle X,Y\rangle X;
\end{aligned}
\]
\[
  \tag{2'}\label{sga3:xx:lemma:2-10:equation:2-prime}
\begin{aligned}
  \operatorname{Ad}(\exp Y)H
    &=H+\bar\alpha(H)Y,&
  \operatorname{Ad}(\exp Y)Y'&=Y',\\
  \operatorname{Ad}(\exp Y)X
    &=X+\langle X,Y\rangle H_{-\alpha}
        -\langle X,Y\rangle Y;
\end{aligned}
\]
\[
\tag{3}\label{sga3:xx:lemma:2-10:equation:3}
\begin{gathered}
  [H,X]=\bar\alpha(H)X,\qquad
  [H,Y]=-\bar\alpha(H)Y,\\
  [X,Y]=\langle X,Y\rangle H_\alpha;
\end{gathered}
\]
\begin{align}
  H_{-\alpha}&=-H_\alpha, \tag{4}\label{sga3:xx:lemma:2-10:equation:4}\\
  \bar\alpha(H_\alpha)&=2. \tag{5}\label{sga3:xx:lemma:2-10:equation:5}
\end{align}

The proof of these formulas is either immediate or an immediate
consequence of formula \textup{(F)} of \SGARefLink{sga3:xx:2.1}{2.1}.

\paragraph{\textbf{Corollary 2.11.}}
\SGARefTarget{sga3:xx:2.11}
\emph{Suppose that \(H_\alpha\) is nonzero on every fiber; by
\textup{\SGARefLink{sga3:xx:lemma:2-10:equation:5}{(5)}}, this holds in particular if \(2\) is invertible on
\(S\). Then}
\[
  X_\alpha\in\Gamma(S,\mathfrak g_\alpha)^\times,\qquad
  X_{-\alpha}\in\Gamma(S,\mathfrak g_{-\alpha})^\times
\]
\emph{are paired if and only if}
\[
  [X_\alpha,X_{-\alpha}]=H_\alpha.
\]

\paragraph{\textbf{2.12.}}
\SGARefTarget{sga3:xx:2.12}
Let \((G,T,\alpha)\) be an \(S\)-elementary system. By \SGARefLink{sga3:xx:1.19}{1.19},
\[
  \operatorname{Centr}(G)=\ker(\alpha)
\]
is a group of multiplicative type and finite type. If \(Q\) is a
subgroup of multiplicative type of \(\operatorname{Centr}(G)\),
then \(G/Q\) is affine over \(S\) \textup{(\SGARefLink{sga3:ix:corollary:2:5}{Exposé IX, 2.5})}, smooth
over \(S\) \textup{(Exposé VI\textsubscript{B}, \SGARefLink{prop:sga3-VIB-9.2}{9.2})}, and has
connected reductive fibers of semisimple rank \(1\)
\textup{(\SGARefLink{sga3:xix:paragraph:1-8}{Exposé XIX, 1.8})}.

Put \(G'=G/Q\). It is a reductive \(S\)-group of semisimple rank
\(1\), and \(T'=T/Q\) is a maximal torus of \(G'\). The open
subset
\[
  U_{-\alpha}\cdot T\cdot U_\alpha
\]
of \(G\) is stable under \(Q\), and its quotient is isomorphic to
\[
  U_{-\alpha}\times_S(T/Q)\times_S U_\alpha.
\]
If \(\alpha'\) denotes the character of \(T'\) induced by
\(\alpha\), the canonical morphism \(G\to G'\) induces
isomorphisms
\[
  \mathfrak g_\alpha\xrightarrow{\sim}\mathfrak g'_{\alpha'},
  \qquad
  \mathfrak g_{-\alpha}\xrightarrow{\sim}\mathfrak g'_{-\alpha'}.
\]
In particular, \(\alpha'\) is a root of \(G'\) relative to \(T'\).
Writing \(\alpha/Q\) for the character
\[
  T/Q\longrightarrow\mathbb G_{m,S}
\]
induced by \(\alpha\), we obtain:

\paragraph{\textbf{Lemma 2.13.}}
\SGARefTarget{sga3:xx:2.13}
\emph{If \(Q\) is a subgroup of multiplicative type of
\(\ker(\alpha)\), then}
\[
  (G/Q,T/Q,\alpha/Q)
\]
\emph{is an elementary system.}
